\documentclass{article}
\usepackage{graphicx} % Required for inserting images

\title{SemEval_task12}
\author{}
\date{January 2026}

\begin{document}

\section{Introduction}
Abductive reasoning the task of inferring the most plausible explanation for a set of incomplete or partially observed facts is a fundamental component of human cognition. Unlike deductive reasoning, which guarantees logically valid conclusions, or inductive reasoning, which generalizes from observed patterns, abductive reasoning operates under uncertainty and requires generating explanatory hypotheses that best account for available evidence. This form of reasoning underlies many natural language understanding tasks, including commonsense inference, narrative completion, multi-hop question answering, and scientific hypothesis formation.

Recent advances in large language models (LLMs) have substantially improved performance across a wide range of reasoning benchmarks. Transformer-based architectures such as the LLaMA family have demonstrated strong capabilities when trained with large-scale pretraining followed by supervised fine-tuning (SFT) and instruction tuning. Despite these advances, prior work has shown that LLMs often struggle with structured explanatory reasoning. In abductive settings, models may produce fluent but logically inconsistent or weakly grounded explanations, reflecting reliance on surface-level correlations rather than principled hypothesis generation.

Most existing approaches rely primarily on supervised fine-tuning, which optimizes next-token likelihood under labeled examples. While effective for distribution alignment, SFT does not explicitly optimize reasoning quality, explanatory coherence, or plausibility under uncertainty. Consequently, models may fail to distinguish between merely plausible continuations and genuinely explanatory hypotheses.

In this work, we investigate whether reinforcement learning (RL) can enhance abductive reasoning in large language models. Reinforcement learning enables direct optimization of task-specific reward signals beyond likelihood maximization. By designing reward functions that capture explanatory plausibility, logical consistency, and coherence with observed evidence, we aim to guide the model toward generating higher-quality abductive inferences.

We propose a reinforcement learning fine-tuning framework built on top of a LLaMA-based instruction-tuned model. Our method augments supervised training with policy optimization using structured reward signals tailored to abductive reasoning tasks. We conduct systematic experiments and ablation studies to evaluate the impact of reward design and training strategy on reasoning performance.

Our contributions are threefold:

We introduce a reinforcement learning framework specifically designed to improve abductive reasoning in large language models.

We analyze the role of reward shaping and optimization strategies through controlled ablation experiments.

We demonstrate consistent improvements over supervised fine-tuning baselines on abductive reasoning benchmarks.

Our findings suggest that reinforcement learning provides a promising direction for moving beyond pattern-based generation toward more structured and explanation-driven reasoning in large language models.

\section{Related Work}

### Related Work

**Li et al. (2023), *Language Models Can Improve Event Prediction by Few‑Shot Abductive Reasoning*** investigates the use of large pretrained language models to generate plausible explanatory hypotheses for event prediction tasks. A key methodological contribution of this work is the use of *few‑shot prompting* to structure model reasoning and improve inference quality. We draw on this insight by similarly structuring model outputs — instead of treating reasoning as implicit — and use explicit reasoning text as intermediate representations. Unlike the few‑shot prompt strategy, our approach trains a verifier model to learn from explicit reasoning errors.

**Zhao et al. (2025), *Advancing Reasoning in Large Language Models*** provides a comprehensive analysis of reasoning capabilities in large language models, categorizing different reasoning paradigms and identifying recurrent failure modes such as shallow pattern matching and lack of contextual grounding. Critically, this work highlights the value of *auxiliary feedback mechanisms and multi‑stage reasoning pipelines* to mitigate reasoning brittleness. We adopt this principle by combining a primary reasoning model (Qwen) with an auxiliary verification model (DeBERTa), enabling systematic detection and correction of reasoning weaknesses.

**Wang et al. (2025), *Causal Event Understanding with Contextual Inference*** (ACL Long Paper 1269) focuses on models that must integrate rich contextual evidence to disambiguate competing causal explanations. This work introduces *contextual document integration and multi‑label causal classification* as core components of its methodology. Inspired by this, our fine‑tuned DeBERTa multi‑label classifier also treats each option as a separate target and encodes the combined event, context, and options to capture nuanced causal signals, using sliding window tokenization to preserve relevant long‑range context.

Together, these studies inform our methodology: we explicitly structure reasoning outputs drawing from few‑shot guidance, incorporate a learned verification stage to address reasoning brittleness, and design multi‑label classification with contextual integration to handle complex abductive inference scenarios.

\section{Method}
\subsection{Base Models}

Meta-Llama-3.1-70B-Instruct

s
microsoft/deberta-v3-base



\subsection{Supervised Fine-Tuning}

3.2.1 ( Fine-Tuned DeBERTa Multi-Label Classification)
For this method, we fine-tuned microsoft/deberta-v3-base on the SemEval Task 12 dataset. Each instance consists of the target event, contextual documents, and four candidate options (A–D). Because the contexts can be long, we applied a sliding window tokenization strategy with overlapping strides to ensure that all relevant context is captured within the model’s maximum sequence length. The gold answers are encoded as a multi-hot vector, allowing the model to predict the correctness of each option simultaneously. The model is trained with a multi-label classification head to minimize binary cross-entropy loss, and the checkpoint with the lowest training loss is saved for evaluation. During inference, the model outputs a vector of probabilities for each option, and a threshold of 0.5 determines predicted correctness.

\subsection{Reinforcement Learning Framework}

\subsection{Reward Design}

\subsection{Qwen Reasoning Generation and DeBERTa Verification}
The system follows a two-stage reasoning and verification pipeline. First, the Qwen model receives the target event, candidate explanations (Option A–D), and contextual documents. Instead of producing only a final prediction, Qwen generates structured outputs for each option in the following format:
Option A is correct: True/False
Reasoning: explanation generated by Qwen
Option B is correct: True/False
Reasoning: explanation generated by Qwen
Option C is correct: True/False
Reasoning: explanation generated by Qwen
Option D is correct: True/False
Reasoning: explanation generated by Qwen
This structured reasoning format allows the model to explicitly justify its abductive decisions and enables downstream verification. The reasoning and correctness labels for all options are then passed to a DeBERTa-based verifier model.
During training, DeBERTa receives the original event, candidate options, contextual documents, Qwen’s True/False predictions, and the reasoning text. These are compared with the golden answers to identify mismatches between Qwen’s predictions and the ground truth. Through this process, DeBERTa learns recurring patterns in Qwen’s errors, such as overgeneralized causal assumptions, weak justification, or reasoning that contradicts contextual evidence. The verifier therefore learns to associate certain reasoning styles with a higher probability of incorrect predictions.
At inference time, Qwen first produces the structured option-level outputs. DeBERTa then evaluates both the correctness labels and the reasoning quality. If the reasoning resembles patterns that previously led to mistakes, DeBERTa flags the prediction and signals that the reasoning is unreliable, effectively acting as a judge. The system can then request Qwen to reconsider the prediction. If repeated reasoning remains inconsistent or unreliable, DeBERTa produces its own final decision based on the learned verification patterns.

\subsection{Zero-shots}
{3.6.1 DeBERTa Zero-Shot Multiple-Choice Classification }

In addition to the reasoning-verification pipeline, we implemented a zero-shot multiple-choice approach using microsoft/deberta-v3-base. Since DeBERTa was not originally trained with a multiple-choice head, each candidate option was treated as a separate “sentence pair” with the target event. Each pair was encoded independently using the model tokenizer, and the output logits were used to assign a score to the option. The option with the highest score was selected as the predicted answer. This approach does not require task-specific fine-tuning and relies solely on the pretrained model representations. The method was applied directly to the structured questions dataset, processing each option independently and comparing the predicted answer to the golden label.


\section{Results}

(Section 3.5: Qwen Reasoning Generation and DeBERTa Verification)
The method described in Section 3.5 performs very well in cases where Qwen’s initial reasoning is largely aligned with contextual evidence, particularly when there is a clear causal relationship between the target event and the candidate options. In such instances, the DeBERTa verifier effectively reinforces correct predictions and filters out minor reasoning inconsistencies. However, the method struggles in cases that require deeper commonsense knowledge or fine-grained accuracy, especially when multiple explanations appear superficially plausible. In these scenarios, the verifier can only partially mitigate errors, as both models may rely on similar implicit assumptions or insufficient context. Overall, the approach demonstrates a clear improvement over the baseline, while highlighting remaining challenges in highly ambiguous or complex reasoning instances.

(3.6.1 DeBERTa Zero-Shot Multiple-Choice Classification )
The zero-shot DeBERTa multiple-choice method achieves approximately 50 percentage  accuracy on the evaluated dataset. While this approach provides a reasonable baseline without task-specific training, performance is limited by the model’s lack of fine-tuning for multi-option reasoning. Errors primarily occur in instances requiring complex abductive inference or context-sensitive justification, where the model cannot reliably distinguish subtle differences between candidate options. Despite these limitations, the method demonstrates that pretrained encoder representations can provide a non-trivial baseline and can serve as a reference point for further improvements through structured reasoning or verification pipelines.

3.2.1 ( Fine-Tuned DeBERTa Multi-Label Classification)
The fine-tuned DeBERTa model achieves approximately 70 percentage accuracy on the development set. This represents a substantial improvement over the zero-shot baseline, demonstrating that task-specific supervision allows the model to correctly identify multiple plausible options while suppressing incorrect selections. The sliding window approach effectively handles the long contextual documents, enabling the model to leverage relevant information across large input spans. Errors mainly occur in instances requiring deep commonsense reasoning or where contextual cues are sparse, but overall the fine-tuned model provides a robust and reliable baseline for multi-label abductive reasoning on Task 12.
\section{Conclusion}




\bibliographystyle{plain}
\bibliography{ref}
\nocite{*}


\end{document}



